\section{Results IV: prospective check}
\label{sec:prospective}

Tier P repeats the information-matched comparison on forecasts made before their outcomes existed. Each issue was sealed by a
combined SHA-256 of its contexts, tasks, issue record, code bundle, run records and all forecast files, computed before the first
day of its scoring window (Table~\ref{tab:seals}). The analysis (secondary family S6) pools both issues and tests TimesFM~3 in
univariate mode against the statistical ensemble, Chronos-2, TimesFM~2.5 and the naive forecast, on MASE and SQL, separately on the
structured domains and on Wikipedia, with series clusters only. With two overlapping issues, Tier P generalises over series, not
over time: it checks the direction of the Tier B results, not their size.

\begin{table}[h]
\centering
\caption{Tier P issues and seals. The scoring window covers target timestamps from two to nine days after the issue day (UTC).}
\label{tab:seals}
\scriptsize
\begin{tabular}{lllrl}
\toprule
Issue day & Sealed (UTC) & Scoring window & Tasks & Combined SHA-256 \\
\midrule
2026-10-30 & 2026-10-30 08:10 & 11-01 to 11-08 & 385 & \texttt{f84a669d33e07ffb\ldots 349843f3} \\
2026-11-02 & pending & 11-04 to 11-11 & -- & -- \\
\bottomrule
\end{tabular}
\end{table}

\textit{[Results to be added after the outcomes of both issues are published (planned 2026-11-10 and 2026-11-13).]}
